%%%%%%%%%%%%%%%%%%%%%%%%%%%%%%%%%%%%%%%%%%%%%%%%%%%%%%%%%%%%%%%%%%%%%%%%%%%
% Función: capítulo de ejemplo de introducción de la tesis por compendio.
%
% Uso: modifica este fichero para adaptar el contenido a tu trabajo.
%%%%%%%%%%%%%%%%%%%%%%%%%%%%%%%%%%%%%%%%%%%%%%%%%%%%%%%%%%%%%%%%%%%%%%%%%%%

\ifthenelse{\equal{\myLanguage}{english}}
{
  \chapter{Introduction}
}
{
  \chapter{Introducción}
}
\label{cha:compendium-introduction}

\ifthenelse{\equal{\myLanguage}{english}}
{
  \section{Context and Motivation}
  \subsection{Research context}
}
{
  \section{Contexto y motivación}
  \subsection{Contexto de la investigación}
}

\ifthenelse{\equal{\myLanguage}{english}}
{
  Use this chapter to introduce the research context, motivation and central problem addressed by the thesis. Explain why the problem is relevant, delimit the scope of the work and provide the reader with the information required to understand the extended summary. This sample also demonstrates the use of acronyms: \gls{AI} and \gls{HMI}.
}
{
  Utiliza este capítulo para presentar el contexto de la investigación, su motivación y el problema central abordado por la tesis. Explica por qué el problema es relevante, delimita el alcance del trabajo y proporciona la información necesaria para comprender el resumen extendido. Este ejemplo también muestra el uso de acrónimos: \gls{AI} y \gls{HMI}.
}

\ifthenelse{\equal{\myLanguage}{english}}
{
  \subsection{Problem statement}

  Define the specific problem addressed by the thesis, its scope and the assumptions that delimit the proposed research.

  \section{Organization of the Thesis}

  Briefly describe the contents of the extended summary and explain how the publications included in Part~\ref{part:compendium-publications} support the thesis narrative.
}
{
  \subsection{Planteamiento del problema}

  Define el problema concreto abordado por la tesis, su alcance y las hipótesis que delimitan la investigación propuesta.

  \section{Organización de la tesis}

  Describe brevemente el contenido del resumen extendido y explica cómo las publicaciones incluidas en la parte~\ref{part:compendium-publications} sustentan el hilo argumental de la tesis.
}

%%% Local Variables:
%%% TeX-master: "../../book"
%%% End:
